\documentclass[10pt,a4paper]{amsart}
\usepackage[margin=19mm]{geometry}
\usepackage{amsmath,amssymb,mathtools}
\usepackage[scr=rsfs,cal=euler]{mathalfa}
\usepackage{xfrac}
\usepackage[unicode,hidelinks,bookmarks=false]{hyperref}
\newcommand{\ie}{i.e.\ }
\newcommand{\cf}{cf.\ }
\newcommand{\resp}{respectively\ }
\newcommand{\Subsection}{\subsection}
\providecommand{\nobreakdash}{\nobreak\hbox{-}}
\setlength{\emergencystretch}{2em}
\multlinegap0pt
\usepackage{bm}
\usepackage{bbm}
\usepackage{dsfont}
\usepackage{xspace}
\usepackage{cancel}
\usepackage{mathtools}
\usepackage{amsthm}
\makeatletter
\@ifundefined{thm}{\newtheorem{thm}{Thm}}{}
\@ifundefined{lemma}{\newtheorem{lemma}[thm]{Lemma}}{}
\makeatother
\DeclareMathOperator{\Ext}{Ext}
\DeclareMathOperator{\Hom}{Hom}
\newcommand{\C}{\mathbb{C}}
\newcommand{\sO}{\mathcal{O}}
\title[Moduli spaces on the Kuznetsov component of Fano threefolds ]{Moduli spaces on the Kuznetsov component of Fano threefolds of index 2}
\author{Matteo Altavilla, Marin Petkovic, Franco Rota}
\date{Source: arXiv:1908.10986v4 (CC BY 4.0)}
\begin{document}
\maketitle

\noindent\emph{Source and licence.} Moduli spaces on the Kuznetsov component of Fano threefolds of index 2. Version arXiv:1908.10986v4 (CC BY 4.0), distributed under the Creative Commons Attribution 4.0 International licence, as recorded in the arXiv licence metadata for this exact version and shown on the abstract page https://arxiv.org/abs/1908.10986v4. Full rights notice: licenses/arxiv-1908.10986-CC-BY-4.0.txt.\par\smallskip

\noindent\textit{Editorial scope.} Counted unit: the Lemma whose source label is \texttt{d+3} in the article (source0.tex lines 961--969, source label \texttt{d+3}), delivered together with the article's own complete original proof of it (source0.tex lines 971--1007, one \texttt{proof} environment of 37 lines). No mathematics is written, repaired or re-derived: statement and proof are the article's own text line for line. Required context retained uncounted: none. Every definition, display and label of those lines is retained unchanged. Omitted from this package and not redistributed: 1--960, 970--970, 1008--2072. Nothing inside a retained excerpt is omitted or altered: every retained body is delivered exactly as the article's lines are written, so no omission receipt of any kind accompanies this package. Citations: the retained excerpts carry no citation command at all, so no bibliography entry of the article is transcribed and no reference is redistributed. Reference adaptation: none was needed and none was made; every reference inside the delivered unit resolves inside it, so no unresolved reference or citation is produced. Printed slips kept literal: no printed word, symbol, hypothesis or number of the counted statement or of its proof was altered. Layout and class: the article's own class is \texttt{amsart} with packages including epigamath, babel, enumerate, tikz-cd, mathtools, amssymb, xy, leftidx, xcolor; the wrapper is the lane's pinned \texttt{amsart} editorial wrapper at 10pt on A4, which restates the theorem-like environments the retained text needs and adds the article's own macro definitions that the retained text needs (\texttt{\textbackslash C}, \texttt{\textbackslash Ext}, \texttt{\textbackslash Hom}, \texttt{\textbackslash sO}), with extra wrapper packages (bm, bbm, dsfont, xspace, cancel, mathtools). The delivered numbering is the unit's own. Assets: no retained span contains a figure environment, a graphics-inclusion command, a PDF-page inclusion, a table or a TikZ picture; no image file is copied into this package, the package declares no asset dependency and no image file is redistributed with it. Licence: version arXiv:1908.10986v4 of this article is distributed under the Creative Commons Attribution 4.0 International licence, as recorded in the arXiv licence metadata for this exact version and shown on the abstract page https://arxiv.org/abs/1908.10986v4. A full rights notice is delivered at licenses/arxiv-1908.10986-CC-BY-4.0.txt. Characters: every retained excerpt is pure ASCII, so the delivered engine, XeLaTeX with its default Latin Modern text font, reports no missing character.
\begin{lemma}\label{d+3}
Let $S\subset Y$ be a hyperplane section, let $p$ a smooth  point of $S$. Then we have
$$\dim \Ext^i(I_{p|S},I_{p|S})=\begin{cases}
0 & \text{ if }i=3\\
2 & \text{ if }i=2\\
d+3 & \text{ if }i=1\\
1  & \text{ if }i=0.\end{cases}
$$  
\end{lemma}

\begin{proof}
Since $I_{p/S}\simeq[\sO_S\rightarrow \C_p]$ in $D^b(Y)$, then 
$R\Hom( I_{p/S},I_{p/S})$ is the totalization of a double complex $K^{\bullet,\bullet}$ isomorphic to 
$$K^\bullet\coloneqq \left[R\Hom(\C_p,\sO_S)\rightarrow R\Hom(\sO_S,\sO_S)\oplus R\Hom(\C_p,\C_p)\rightarrow R\Hom(\sO_S,\C_p)\right],$$ so we may compute $\Ext^i( I_{p/S},I_{p/S})$ using a spectral sequence
$$E_1^{p,q}=H^q(K^{\bullet,p})\Rightarrow H^{p+q}(K^\bullet).$$
Then the first page $E^{p,q}_1$ of the spectral sequence is

\begin{center}
%$E_2^{p,q}$\quad = \quad
\begin{tabular}{c|cc}
$\vdots$ & $\vdots$ & $\vdots$\\
 $\Ext^1(\C_p,\sO_S)$ & $\Ext^1(\sO_S,\sO_S)\oplus \Ext^1(\C_p,\C_p)$ &$\Ext^1(\sO_S,\C_p)$\\
$\Hom(\C_p,\sO_S)$ & $\Hom(\sO_S,\sO_S)\oplus \Hom(\C_p,\C_p)$ &$\Hom(\sO_S,\C_p)$\\
\hline
($p=-1$)&($p=0$)&($p=1$)
\end{tabular} 
\end{center}
with arrows pointing to the right. The dimensions of the vector spaces above are given in the table below.
\begin{center}
\begin{tabular}{c|cc}
1 & 1 &0\\
1 & 3 &0\\
0 & $d+4$ &1\\
0 & 2 &1\\
\hline
\end{tabular}
\end{center}
The map in the top row is an isomorphism $\Ext^3(\C_p,\sO_S)\rightarrow \Ext^3(\C_p,\C_p)$. The map in the bottom row is  surjective because $\Hom(\sO_S,\sO_S)\rightarrow \Hom(\sO_S,\C_p)$ is.

Consider the maps in the third and second rows. We claim that if $p$ is smooth in $S$, then these two are both non-zero, and if $p$ is singular, then they are both $0$. 
For example, the map in the third row is $\Ext^1(\C_p,\C_p)\rightarrow \Ext^1(\sO_S,\C_p)$. By applying the functor $\Hom(-,\C_p)$ to the sequence $I_{p|S} \to \sO_S \to\C_p$, one sees that its kernel is $\Hom(I_{p|S},\C_p)$. On the other hand we have 
\[ \Hom(I_{p|S},\C_p)\simeq \Hom_S(I_{p|S},\C_p)\simeq \Ext^1_S(\C_p,\C_p)\simeq T_p S,\]
and the latter has rank 2, resp. 3, if $p$ is smooth, resp. singular, in $S$.
The map in the second row is $\Ext^2(\C_p,\sO_S)\to \Ext^2(\C_p,\C_p)$, and the argument for the claim is similar.

This shows that the second page of the spectral sequence is supported on the central column of the table, and hence $E_2^{p,q}$ abuts to the claimed dimensions.
\end{proof}

\end{document}
